\section{Models, observations and coordination}\label{sec:01-introduction-1}
A coordinator can make the right decision without identifying a whole utility
table; identifying a relabeling need not reveal enough to optimize a task.
Similarly, a symmetric mixed equilibrium can exist while an efficient
realization requires a coordination resource. We examine these distinctions
in finite models with explicit timing, information and receipts.

Other-play and label-free coordination already investigate symmetry and
independently constructed policies \citep{hu2020,treutlein2021}. Neural
solvers and equivariant approximators directly study equilibrium selection
and welfare limitations \citep{marris2022,duan2023}. Tewolde et al.\
\citep{tewolde2025} connect game automorphisms with computation and recall
the classical existence of invariant mixed equilibria. The contribution
assessed here is therefore a set of precise model comparisons, with novelty
unresolved, rather than a new general symmetry or equilibrium theory.

The mathematical comparisons concern finite coordination models with explicit
information and implementation resources. The action-orbit computation results
use typed game automorphisms; the utility comparisons retain a specified common
receipt scale. Full structural classification and general assignment fairness
are separate questions and are not prerequisites for the results here.

All initial games and state sets are finite. A typed relabeling moves players
and bijects their action sets; products act right to left. Exact maps preserve
numerical utilities, cardinal maps preserve each recipient's expected-utility
preferences, and ordinal maps preserve pure-outcome preferences. Interpersonal
sums below use explicitly calibrated receipts. For software agents such
objectives are design specifications, not independently established welfare
interests. We distinguish pure equilibria, independent mixed profiles,
lotteries over equilibria, and private recommendations.

Sections~\ref{sec:02-information-1}--\ref{sec:05-state-1} give the allocation,
equilibrium-selection, private-cost and state models. The main refinements are
the guaranteed-success screening bound
(Proposition~\ref{result:07-phase11-1}), the sharp trusted-query frontier
(Proposition~\ref{result:07-phase11-2}), and the explicitly tabulated local-factor
reduction (Proposition~\ref{result:07-phase11-3}). The representation boundary
and implementation costs explain why these conclusions require different
resources. Each result retains its finite model and stated assumptions; no
claim about general-purpose language-model behaviour follows.
